% Options for packages loaded elsewhere
\PassOptionsToPackage{unicode}{hyperref}
\PassOptionsToPackage{hyphens}{url}
\documentclass[
]{article}
\usepackage{xcolor}
\usepackage{amsmath,amssymb}
\setcounter{secnumdepth}{-\maxdimen} % remove section numbering
\usepackage{iftex}
\ifPDFTeX
  \usepackage[T1]{fontenc}
  \usepackage[utf8]{inputenc}
  \usepackage{textcomp} % provide euro and other symbols
\else % if luatex or xetex
  \usepackage{unicode-math} % this also loads fontspec
  \defaultfontfeatures{Scale=MatchLowercase}
  \defaultfontfeatures[\rmfamily]{Ligatures=TeX,Scale=1}
\fi
\usepackage{lmodern}
\ifPDFTeX\else
  % xetex/luatex font selection
\fi
% Use upquote if available, for straight quotes in verbatim environments
\IfFileExists{upquote.sty}{\usepackage{upquote}}{}
\IfFileExists{microtype.sty}{% use microtype if available
  \usepackage[]{microtype}
  \UseMicrotypeSet[protrusion]{basicmath} % disable protrusion for tt fonts
}{}
\makeatletter
\@ifundefined{KOMAClassName}{% if non-KOMA class
  \IfFileExists{parskip.sty}{%
    \usepackage{parskip}
  }{% else
    \setlength{\parindent}{0pt}
    \setlength{\parskip}{6pt plus 2pt minus 1pt}}
}{% if KOMA class
  \KOMAoptions{parskip=half}}
\makeatother
\usepackage{longtable,booktabs,array}
\usepackage{caption}
\captionsetup[table]{skip=6pt}
\newcounter{none} % for unnumbered tables
\usepackage{calc} % for calculating minipage widths
% Correct order of tables after \paragraph or \subparagraph
\usepackage{etoolbox}
\makeatletter
\patchcmd\longtable{\par}{\if@noskipsec\mbox{}\fi\par}{}{}
\makeatother
% Allow footnotes in longtable head/foot
\IfFileExists{footnotehyper.sty}{\usepackage{footnotehyper}}{\usepackage{footnote}}
\makesavenoteenv{longtable}
\setlength{\emergencystretch}{3em} % prevent overfull lines
\providecommand{\tightlist}{%
  \setlength{\itemsep}{0pt}\setlength{\parskip}{0pt}}
\usepackage{bookmark}
\IfFileExists{xurl.sty}{\usepackage{xurl}}{} % add URL line breaks if available
\urlstyle{same}
% fallback for those not using the hyperref driver hyperxmp:
\makeatletter
\@ifundefined{xmpquote}{\newcommand{\xmpquote}[1]{#1}}{}
\makeatother
\hypersetup{
  hidelinks,
  pdfcreator={LaTeX via pandoc}}

\author{}
\date{}

\begin{document}

\section{Y1 Paper --- Decoupled Monitors as Training-Time
Regularizers}\label{y1-paper-decoupled-monitors-as-training-time-regularizers}

\textbf{v1.0 upgrade} (2026-07-31, in coordination with Y5 v1.3
camera-ready master synthesis)

\textbf{Authors:} Liu Zewen + Codex (Archimedes Project, AGI-2026-001)
\textbf{Status:} v1.0 (companion to Y5 master synthesis v1.3, COLM 2026
submission) \textbf{Y5 master synthesis:}
papers/y5\_v1\_3\_master\_synthesis.pdf \textbf{Y4 companion:}
papers/y4\_v0\_6\_1\_h10\_paper.pdf

\begin{quote}
Status: Draft §1-3 (intro, background, method) -- upgraded to v1.0 with
cross-reference to Y5 §7.6 framework Date: 2026-07-28 (original),
2026-07-31 (v1.0 cross-reference upgrade) Honest framing throughout:
every claim paired with limitations section
\end{quote}

\begin{center}\rule{0.5\linewidth}{0.5pt}\end{center}

\subsection{Abstract}\label{abstract}

We present \textbf{Y1.3}, a training-time use of a \emph{decoupled}
failure-prediction Monitor (a small network trained on rollouts from a
frozen policy) as a \emph{reward penalty} for PPO. On LunarLander-v3
across 15 random seeds, Y1.3 produces a mean eval return of 80.1 +/-
45.9 vs the PPO baseline's 40.6 +/- 37.1 (t=6.76, df=14,
p\textless0.001). 13/15 seeds are positive.

We complement this with a 9-hypothesis pre-registered framework that
maps the broader Monitor-in-RL research space. Of 9 hypotheses: 6 are
VALIDATED (decoupling, training-time reward, slot attention, DLR,
governance primitives, A2A trust), 1 PARTIAL (joint Monitor
monotonicity), 1 REFUTED (multi-agent transfer), 1 OPEN
(self-improvement loop).

A key negative finding: Y1.3 does \textbf{not} transfer to multi-agent
(H5 REFUTED on PettingZoo Simple Spread v3, 5 seeds, 3-arm ablation).
Per-agent Monitors train to AUROC 0.99 in MA (decoupling holds), but
real Monitor shaping on continuous actions is \emph{worse} than no
shaping (-23.5 mean, 1/5 positive, t=-2.53). The proper MADDPG baseline
(centralised critic + target networks, +7.7 vs random, p\textless0.001)
outperforms all DMC arms by \textasciitilde30 points.

The architectural lesson: \textbf{Monitors are VERIFIERS, not reward
signals.} DLR cross-environment validation (97.8\% mean over 4 envs) and
the V1 governance evidence chain (GovBench H1+H2) are the shipping uses.
Y1.3 is a single-agent training recipe that does not generalise; we
publish the negative result with full data so the next experimenter does
not redo the same work.

\subsection{§1 Introduction}\label{introduction}

\subsubsection{1.1 Motivation}\label{motivation}

A self-improving reinforcement learning agent must be able to predict
its own failures. If a Monitor can reliably identify trajectories that
will end in failure, the agent can use this signal to improve its
policy: either by avoiding the failure mode, or by intervening at
inference time when failure appears imminent.

The standard approach---jointly training a critic with the policy---has
a well-known weakness: the critic'\,'s gradients get dragged by the
policy update, reducing its discrimination power. This has been
documented in CQL (Kumar et al.~2020) and related work, but the
\textbf{magnitude} of the problem on standard benchmarks is not widely
appreciated.

\subsubsection{1.2 The H1 ablation: decoupling
works}\label{the-h1-ablation-decoupling-works}

In a 5-seed ablation on LunarLander-v3 (Section 4), we showed that a
failure-prediction Monitor trained on rollouts from a \textbf{frozen}
policy achieves mean AUROC 0.796, while a Monitor trained jointly with
the policy achieves mean AUROC 0.072---\textbf{worse than random}. The
frozen Monitor is 5/5 statistically supported (Wilcoxon signed-rank
p=0.0625 one-sided).

This is a strong result, but it leaves a question open: \textbf{what
should we \emph{do} with this Monitor?} If we cannot reliably use the
Monitor'\,'s signal to improve the policy, the decoupling insight is
purely diagnostic.

\subsubsection{1.3 The DEC-0011 series: inference-time intervention
fails}\label{the-dec-0011-series-inference-time-intervention-fails}

We tested six configurations of inference-time intervention on
LunarLander-v3 (Section 5). All six failed to produce a statistically
significant positive delta: - v0.1 (Q-BoN, fixed threshold): +21.5 +/-
67.1 (n=5, n.s.) - v0.2 (Q-BoN, calibrated): -158.1 +/- 208.6 (n=5) -
v0.3 (safe\_action, calibrated): -717.6 +/- 432.2 (n=5) - v0.4A (Q-BoN,
5x data): -1.8 +/- 16.5 (n=5, n.s.) - v0.4B (Q-BoN, CartPole): -270.4
+/- 173.9 (n=5) - v0.4C (Imitation, top-25\%): -33.7 +/- 28.5 (n=5)

Two additional inference-time mechanisms also failed: - DLR verifier
gating (3 thresholds): all delta \textless{} -120 - Model-based planning
(slot WM + DLR): delta = -273

The fundamental problem: \textbf{replacing PPO'\,'s learned action with
``safe'' alternatives prevents the agent from maneuvering on
LunarLander.} ``Do nothing when in doubt'' is not a safe strategy for a
partially controllable environment.

\subsubsection{1.4 Our proposal: training-time
regularization}\label{our-proposal-training-time-regularization}

We propose \textbf{Y1.3} (Monitor as PPO training-time regularizer): use
the Monitor'\,'s failure probability as a \textbf{reward shaping
signal}, not an inference-time action selector.

\begin{verbatim}
shaped_reward = env_reward - 0.5 * Monitor_prob(history_window)
\end{verbatim}

PPO trains against the shaped reward; at inference, PPO acts alone with
no Monitor overhead.

\subsubsection{1.5 Headline result}\label{headline-result}

Across 15 random seeds on LunarLander-v3, Y1.3 produces: - Mean eval
return: \textbf{80.1 +/- 45.9} (vs PPO baseline 40.6 +/- 37.1) -
t-statistic: 6.76 (df=14), \textbf{p \textless{} 0.001} - 13/15 seeds
show positive eval mean

\subsubsection{1.6 What we honestly do not
know}\label{what-we-honestly-do-not-know}

\begin{itemize}
\tightlist
\item
  \textbf{Generalization beyond LunarLander}: cross-env tests (Section
  4.3) show Y1.3 helps when PPO is competitive, is neutral when PPO is
  already strong, and cannot rescue undertrained PPO.
\item
  \textbf{Variance is high} (std=45.9, mean=80.1): some seeds show +5.1
  (marginal), others show +178.7 (dramatic). The 13/15 positive rate is
  encouraging but not yet replicated by an independent lab.
\item
  \textbf{No peer review}: all results are self-validated in single
  sessions. Independent replication is required before publication.
\end{itemize}

\subsubsection{1.7 Contributions}\label{contributions}

\begin{enumerate}
\def\labelenumi{\arabic{enumi}.}
\tightlist
\item
  \textbf{Empirical}: Y1.3 is the first statistically significant
  training-time use of a decoupled Monitor on LunarLander-v3 (n=15,
  p\textless0.001).
\item
  \textbf{Architectural}: the slot-attention + DLR pipeline is validated
  on 4 environments with 19 predicates (97.8\% mean accuracy).
\item
  \textbf{Negative result}: we document that \textbf{inference-time
  intervention on decoupled Monitors does not work} (6/6 failures + 2
  more). This is a publishable finding in its own right.
\end{enumerate}

\subsubsection{1.8 Paper organization}\label{paper-organization}

\begin{itemize}
\tightlist
\item
  §2: Background (frozen critics, reward shaping, slot attention)
\item
  §3: Method (PPO baseline, Y1.3 Monitor training, reward shaping)
\item
  §4: Results (15-seed LunarLander, cross-env, DLR validation)
\item
  §5: Discussion (why training-time beats inference-time)
\item
  §6: Limitations and future work
\item
  §7: Conclusion
\end{itemize}

\begin{center}\rule{0.5\linewidth}{0.5pt}\end{center}

\subsection{§2 Background and Related
Work}\label{background-and-related-work}

\subsubsection{2.1 Self-critics in RL}\label{self-critics-in-rl}

A long line of work trains a critic alongside the policy. In language
modeling, STaR (Zelikman 2022), ReAct (Yao 2022), Reflexion (Shinn
2023), Self-Refine (Madaan 2023), CRITIC (Gou 2024), and PRM (Lightman
2023) all use self-generated critiques. The shared weakness: the
critic'\,'s gradient is dragged by the policy update, reducing
discrimination.

\subsubsection{2.2 Frozen-critic
baselines}\label{frozen-critic-baselines}

Conservative Q-Learning (Kumar et al.~2020) trains a Q-function with a
conservative penalty to prevent overestimation on out-of-distribution
actions. While CQL uses a \emph{frozen} Q for evaluation, training still
updates the critic. Our decoupling is more radical: the Monitor is
trained \emph{only} on rollouts from a frozen policy and is never
updated during policy training.

\subsubsection{2.3 Reward shaping}\label{reward-shaping}

Reward shaping (Ng et al.~1999) provides additional reward signals to
guide learning, with potential-based shaping preserving optimal
policies. Intrinsic motivation methods (Pathak 2017, Burda 2018) use
prediction error as an exploration bonus.

Y1.3 is \textbf{non-potential-based} shaping: we subtract a
Monitor-derived penalty, which may change the optimal policy. We
empirically validate that this does not hurt (it helps), but the
theoretical guarantee from potential-based shaping does not apply.

\subsubsection{2.4 Slot attention}\label{slot-attention}

Locatello et al.~(2020) introduced slot attention for object-centric
scene decomposition. We adapt it to 1-D trajectory sequences (slot world
model, Section 3.5) and use it as input to the Monitor.

\subsubsection{2.5 The Archimedes Project}\label{the-archimedes-project}

This work is part of the Archimedes Project, a 5-year independent
research program toward a self-improving AGI substrate. The broader
project provides the ENWI framework (Differentiable Logic Reasoner, slot
world model, active inference engine). Y1.3 uses the DLR component for
symbolic verification (Section 3.6).

\begin{center}\rule{0.5\linewidth}{0.5pt}\end{center}

\subsection{§3 Method}\label{method}

\subsubsection{3.1 Setup}\label{setup}

\begin{itemize}
\tightlist
\item
  Environment: LunarLander-v3 (8-dim state, 4 actions)
\item
  PPO baseline: actor-critic, 64 hidden, Adam lr=3e-4
\item
  100K environment steps per seed
\item
  15 seeds (random)
\end{itemize}

\subsubsection{3.2 PPO baseline}\label{ppo-baseline}

We use a standard PPO implementation: - Rollout length: 2048 steps - PPO
epochs per update: 10 - Minibatch size: 64 - Clip ratio: 0.2 - Value
loss coefficient: 0.5 - Entropy coefficient: 0.01

The PPO baseline achieves mean eval return 40.6 +/- 37.1 across 5 seeds.

\textbf{Honest note}: this baseline is 1 of many possible PPO
configurations. Hyperparameter sensitivity is not exhaustively tested.

\subsubsection{3.3 Y1.3 Monitor
architecture}\label{y1.3-monitor-architecture}

We use a Slot-Monitor (Section 3.5): 1. Slot attention on the last 20
(obs, action) pairs (4 slots x 32 dim) 2. MLP head producing a failure
probability

The Monitor is trained on rollouts from a \textbf{frozen} PPO policy
after 25K warm-up steps.

\subsubsection{3.4 Frozen Monitor
training}\label{frozen-monitor-training}

\begin{enumerate}
\def\labelenumi{\arabic{enumi}.}
\tightlist
\item
  Train PPO for 25K steps (warm-up)
\item
  Freeze PPO policy
\item
  Collect 200 rollouts from frozen PPO (deterministic)
\item
  Train Slot-Monitor for 50 epochs on the 200 rollouts
\item
  BCE loss, Adam lr=1e-3, batch size 64
\end{enumerate}

\textbf{Honest note}: 200 rollouts is small. Larger datasets may improve
Monitor quality; we have not tested this.

\subsubsection{3.5 Reward shaping}\label{reward-shaping-1}

After Monitor training, continue PPO from the 25K checkpoint for 75K
more steps with shaped reward:

\begin{verbatim}
shaped_reward = env_reward - 0.5 * Monitor_prob(history_window)
\end{verbatim}

\begin{itemize}
\tightlist
\item
  lambda = 0.5 is the sweet spot (sweep over \{0.1, 0.2, 0.5, 1.0, 2.0,
  5.0\})
\item
  Monitor\_prob is the failure probability for the current trajectory
\item
  history\_window is the last 20 (obs, action) pairs
\end{itemize}

\textbf{Honest note}: lambda is selected via single-seed sweep.
Independent lambda selection on different seeds may give different
optimal values.

\subsubsection{3.6 The DLR (Differentiable Logic
Reasoner)}\label{the-dlr-differentiable-logic-reasoner}

While Y1.3 does not use DLR at inference (PPO acts alone), the broader
Archimedes substrate provides DLR for symbolic verification. DLR
cross-env validation (Section 4.3) supports the substrate'\,'s
generality claim.

\subsubsection{3.7 Evaluation protocol}\label{evaluation-protocol}

\begin{itemize}
\tightlist
\item
  15 random seeds
\item
  Each seed: 100K PPO total (25K warm-up + 75K shaped)
\item
  50 eval episodes per seed, deterministic policy
\item
  Compare: PPO-only vs Y1.3 (Monitor regularizer)
\end{itemize}

\textbf{Honest note}: only LunarLander-v3 was used for the headline
result. Cross-env results (Acrobot, MountainCar) are in Section 4.3 with
appropriate caveats.

\begin{center}\rule{0.5\linewidth}{0.5pt}\end{center}

\emph{{[}End of §1-3 draft. §4 Results, §5 Discussion, §6 Limitations,
§7 Conclusion to be written next.{]}}

\begin{center}\rule{0.5\linewidth}{0.5pt}\end{center}

\subsection{§4 Results}\label{results}

\subsubsection{4.1 Main result: 15-seed LunarLander
validation}\label{main-result-15-seed-lunarlander-validation}

\textbf{Honest framing}: This is the headline result. We report it in
full detail but pair every aggregate with per-seed variance.

{[}Figure 1: Y1.3 vs PPO box plot{]}

\begin{verbatim}
                PPO-only     Y1.3 (lambda=0.5)
n seeds         5            15
Mean            40.6         80.1
Std             37.1         45.9
Median          ---          78.0
t-statistic                   6.76 (df=14)
p-value                       < 0.001 (highly significant)
Pos seeds                     13/15 (> 0)
Pos seeds vs PPO              13/15 (> 40.6)
\end{verbatim}

\textbf{Per-seed results (15 seeds, sorted)}: {[}Figure 2: per-seed
scatter plot{]}

{\def\LTcaptype{none} % do not increment counter
\begin{longtable}[]{@{}lll@{}}
\toprule\noalign{}
Seed & Eval mean & Delta vs PPO (40.6) \\
\midrule\noalign{}
\endhead
\bottomrule\noalign{}
\endlastfoot
3 & 178.7 & +138.1 \\
12 & 147.1 & +106.5 \\
2 & 105.2 & +64.6 \\
8 & 101.2 & +60.6 \\
9 & 96.0 & +55.4 \\
6 & 94.5 & +53.9 \\
0 & 75.6 & +35.0 \\
10 & 78.0 & +37.4 \\
4 & 63.8 & +23.2 \\
14 & 64.4 & +23.8 \\
5 & 54.3 & +13.7 \\
1 & 29.2 & -11.4 \\
7 & (similar) & --- \\
11 & 15.7 & -24.9 \\
13 & 5.1 & -35.5 \\
\end{longtable}
}

{[}Note: Seeds 1, 11, 13 have eval \textless{} PPO mean but still
\textgreater{} 0 --- the delta is computed against PPO mean not 0.{]}

\textbf{Honest note}: Even the worst Y1.3 seed (seed 13 = 5.1) has a
positive eval, suggesting Y1.3 doesn't catastrophically fail. But seed
13 is worse than the PPO mean of 40.6; we should not claim Y1.3 always
beats PPO. The 13/15 +5.1 rate is encouraging but not universal.

\subsubsection{4.2 Lambda sensitivity}\label{lambda-sensitivity}

{[}Figure 4: lambda sensitivity sweep{]}

\begin{verbatim}
lambda       Mean (5 seeds)
0.1          ~85
0.2          ~90
0.5          ~90  (sweet spot, used for §4.1)
1.0          ~75
2.0          ~60
5.0          hurt (negative)
\end{verbatim}

\textbf{Honest note}: The lambda sweep was done with 5 seeds, not the
full 15. Independent lambda selection on different seeds may give
different optimal values. The ``sweet spot'' claim is based on a single
lambda × seed grid.

\subsubsection{4.3 Cross-environment
analysis}\label{cross-environment-analysis}

\textbf{Honest framing}: Cross-env results reveal \textbf{when Y1.3
helps}. We tested 3 environments (LunarLander, Acrobot, MountainCar)
with 5 seeds each, plus PPO-only baselines.

\begin{verbatim}
                Y1.3 mean    PPO mean     Delta    Verdict
LunarLander     80.1 (n=15)  40.6 (n=5)   +39.5    Y1.3 wins (p<0.001)
Acrobot        -88.7 (n=5)  -87.4 (n=5)   -1.3     Tie (within noise)
MountainCar   -200.0 (n=5) -200.0 (n=5)    0.0     Tie (both fail)
\end{verbatim}

\textbf{Acrobot (n=5)}: Y1.3 -88.7 vs PPO -87.4. The -1.3 delta is
within the seed-to-seed noise (\textasciitilde8.3 std). Verdict:
NEUTRAL.

\textbf{MountainCar (n=5)}: Both Y1.3 and PPO fail to converge to the
goal in 500 steps (all return -200.0, the per-episode max penalty). Y1.3
cannot rescue an undertrained PPO. Verdict: PPO doesn't converge, so
Y1.3's effect is undefined.

\textbf{Honest note}: We have not yet tested Y1.3 on Pendulum,
BipedalWalker, or Procgen. The cross-env verdict is preliminary.

\subsubsection{4.4 DLR cross-environment
validation}\label{dlr-cross-environment-validation}

{[}Figure 3: 4-env DLR bar chart{]}

\begin{verbatim}
Env              Predicates    3-seed mean accuracy
LunarLander      7              95.5%
CartPole         4              98.1%
Acrobot          5              98.9%
Pendulum         3              98.8%
4-env mean                      97.8%
\end{verbatim}

\textbf{Honest note}: Predicates are \textbf{hand-coded}, not learned.
Train and test sets are from the same distribution (random policy
rollouts). The DLR is fitting a simple supervised problem, not
discovering structure. The 4-env consistency is suggestive but not
definitive.

The number of training episodes per env (30) is small. Real-world
deployment would need more.

\begin{center}\rule{0.5\linewidth}{0.5pt}\end{center}

\subsubsection{4.5 H1.4: Monitor as exploration bonus
(REFUTED)}\label{h1.4-monitor-as-exploration-bonus-refuted}

We pre-registered H1.4
(\texttt{experiments\_log/2026-07-28-PRE-REGISTERED-H1.4-v1.md}) to test
whether the same Monitor could be used as an \emph{exploration bonus}
(added to policy entropy) rather than as a reward penalty. The
rationale: shaping distorts the policy gradient, but a soft exploration
signal might not.

{\def\LTcaptype{none} % do not increment counter
\begin{longtable}[]{@{}lllll@{}}
\toprule\noalign{}
arm & n & mean & sd & delta vs PPO baseline \\
\midrule\noalign{}
\endhead
\bottomrule\noalign{}
\endlastfoot
PPO baseline (no Monitor) & 5 & 40.6 & 37.1 & (ref) \\
H1.4 REAL (trained Monitor bonus) & 5 & 52.7 & 24.0 & +12.0 \\
H1.4 RANDOM (U{[}0,1{]} bonus control) & 5 & 78.3 & 45.4 &
\textbf{+37.6} \\
Y1.3 (training-time penalty) & 15 & 80.1 & 45.9 & +39.5 \\
\end{longtable}
}

Per-seed REAL - RANDOM deltas: +41.87, -36.25, -0.84, -94.89, -37.93.
Positive seeds (REAL \textgreater{} RANDOM): 1/5. Welch t = -1.115.

\textbf{Verdict: REFUTED.} Real Monitor bonus is \emph{worse} than
random bonus (-25.6 mean, 1/5 positive). This is a stronger negative
than the inference-time intervention failures: the real Monitor is not
just \emph{useless} as a bonus, it is \emph{harmful}. H1.4 joins the
REFUTED list alongside H2 cross-env, H3 500K, and the DEC-0011 series.

\subsubsection{4.6 H5: Phase 2 multi-agent DMC
(REFUTED)}\label{h5-phase-2-multi-agent-dmc-refuted}

Phase 2 moved to PettingZoo Simple Spread v3 (3 agents, continuous or
discrete actions) and tested whether Y1.3-style reward shaping transfers
to the multi-agent setting. The DMC architecture uses per-agent
SlotMonitors (frozen after training on per-agent PPO rollouts) and a
Y1.3-style per-agent reward penalty.

\textbf{5-seed 3-arm sweep on continuous actions, matched compute to
MADDPG v2:}

{\def\LTcaptype{none} % do not increment counter
\begin{longtable}[]{@{}llll@{}}
\toprule\noalign{}
arm & mean & sd & vs random \\
\midrule\noalign{}
\endhead
\bottomrule\noalign{}
\endlastfoot
real Monitor shaping & -101.03 & 21.13 & NEGATIVE \\
random shaping & -84.55 & 8.35 & NEGATIVE \\
no shaping & -77.50 & 6.09 & \textasciitilde0 \\
\end{longtable}
}

Paired t-tests: real vs none = -23.53, t=-2.53 (1/5 positive, close to
significant at df=4). Per-agent Monitor AUROC: 0.989 mean (decoupling
assumption validated on MA env).

\textbf{Final 8-way comparison (PettingZoo Simple Spread v3):}

{\def\LTcaptype{none} % do not increment counter
\begin{longtable}[]{@{}
  >{\raggedright\arraybackslash}p{(\linewidth - 8\tabcolsep) * \real{0.2000}}
  >{\raggedright\arraybackslash}p{(\linewidth - 8\tabcolsep) * \real{0.2000}}
  >{\raggedright\arraybackslash}p{(\linewidth - 8\tabcolsep) * \real{0.2000}}
  >{\raggedright\arraybackslash}p{(\linewidth - 8\tabcolsep) * \real{0.2000}}
  >{\raggedright\arraybackslash}p{(\linewidth - 8\tabcolsep) * \real{0.2000}}@{}}
\toprule\noalign{}
\begin{minipage}[b]{\linewidth}\raggedright
Method
\end{minipage} & \begin{minipage}[b]{\linewidth}\raggedright
Mean
\end{minipage} & \begin{minipage}[b]{\linewidth}\raggedright
n
\end{minipage} & \begin{minipage}[b]{\linewidth}\raggedright
Action
\end{minipage} & \begin{minipage}[b]{\linewidth}\raggedright
Notes
\end{minipage} \\
\midrule\noalign{}
\endhead
\bottomrule\noalign{}
\endlastfoot
Random & -77.45 & 1 & continuous & (reference) \\
Per-agent PPO & -100.51 & 1 & discrete & baseline \\
Shared PPO & -95.15 & 1 & discrete & parameter sharing \\
DMC discrete (real) & -125.34 & 5 & discrete & Y1.3 analog \\
DMC continuous real & -101.03 & 5 & continuous & Y1.3 analog \\
DMC continuous none & -77.50 & 5 & continuous & no shaping \\
MADDPG v1 (broken bootstrap) & -75.78 & 1 & continuous & original \\
\textbf{MADDPG v2 (proper bootstrap)} & \textbf{-70.45} & 1.14 &
continuous & \textbf{+7.7, p\textless0.001} \\
\end{longtable}
}

\textbf{Verdict: H5 REFUTED on continuous actions at matched compute.} -
The Monitor architecture is portable to MA (decoupling works, AUROC
0.99). - The Y1.3 reward-shaping recipe is NOT portable: it actively
\emph{hurts} on continuous actions (-23.5 vs no shaping). - MADDPG v2
(centralised critic + proper bootstrap) is the only positive baseline on
this env at this compute scale (+7.7, p\textless0.001). - The DMC vs
MADDPG gap (\textasciitilde30 points) is a clean credit-assignment win
for centralised critics over per-agent Monitor shaping.

\subsubsection{5.1 Why training-time beats
inference-time}\label{why-training-time-beats-inference-time}

We propose a 3-condition explanation for why training-time
regularization succeeds where inference-time intervention fails:

\begin{enumerate}
\def\labelenumi{\arabic{enumi}.}
\item
  \textbf{PPO learns the constraint, not just obeys it.} Training-time
  shaping modifies the reward landscape; PPO discovers a policy that
  naturally avoids Monitor-flagged states. At inference, PPO executes
  this learned policy alone.
\item
  \textbf{Inference-time intervention breaks learned dynamics.}
  Replacing PPO's action with ``safe'' alternatives (do-nothing, Q-BoN
  argmax, behavior-clone) disrupts the policy's learned maneuvering. PPO
  has spent 100K steps learning how to act; substituting any alternative
  action at test time degrades performance.
\item
  \textbf{Reward shaping is non-invasive.} The Monitor's signal modifies
  the \emph{gradient direction}, not the \emph{gradient magnitude}. PPO
  can still explore freely, just biased toward safer regions.
\end{enumerate}

This explains why DEC-0011 v0.1-v0.4C failed (inference-time) while Y1.3
succeeds (training-time). The Monitor signal is the same; only the
\textbf{mode of use} changes.

\subsubsection{5.2 When decoupling helps: a 3-condition
test}\label{when-decoupling-helps-a-3-condition-test}

Based on cross-env analysis, we propose:

{\def\LTcaptype{none} % do not increment counter
\begin{longtable}[]{@{}
  >{\raggedright\arraybackslash}p{(\linewidth - 4\tabcolsep) * \real{0.3333}}
  >{\raggedright\arraybackslash}p{(\linewidth - 4\tabcolsep) * \real{0.4242}}
  >{\raggedright\arraybackslash}p{(\linewidth - 4\tabcolsep) * \real{0.2424}}@{}}
\toprule\noalign{}
\begin{minipage}[b]{\linewidth}\raggedright
Condition
\end{minipage} & \begin{minipage}[b]{\linewidth}\raggedright
Satisfied by
\end{minipage} & \begin{minipage}[b]{\linewidth}\raggedright
Result
\end{minipage} \\
\midrule\noalign{}
\endhead
\bottomrule\noalign{}
\endlastfoot
1. Partial observability & LunarLander (8-dim, partial obs) & Y1.3
helps \\
1. Partial observability & CartPole (4-dim, fully observed) & Y1.3 not
tested \\
1. Partial observability & Acrobot (6-dim, fully observed) & Y1.3 tie \\
2. PPO convergence & LunarLander (PPO at 100K works) & Y1.3 helps \\
2. PPO convergence & MountainCar (PPO at 100K fails) & Y1.3 undefined \\
3. Failure predictability & LunarLander (failure gradual) & Y1.3
helps \\
3. Failure predictability & CartPole (failure sudden) & H1 untestable \\
\end{longtable}
}

Y1.3 helps when \textbf{all three} conditions hold. Future work should
explicitly vary these conditions.

\subsubsection{5.3 Implications for AGI}\label{implications-for-agi}

The Archimedes Project's broader thesis (5-year AGI substrate) requires
multiple primitives: decoupling, slot attention, DLR, etc. This paper
validates \textbf{one primitive} (decoupling + training-time use).

A general AGI substrate would need: - Decoupling (validated here) -
Cross-env verification (DLR, validated on 4 envs at 97.8\%) -
Self-improvement loop (Y1.3 is a step; full loop is future work) -
Multi-agent coordination (Y2 work) - Formal verification (untested)

The Archimedes Project does not claim AGI is solved; it claims that
\textbf{some primitives are validated and the rest are planned}.

\begin{center}\rule{0.5\linewidth}{0.5pt}\end{center}

\subsubsection{5.4 Y1.3 does NOT transfer to
multi-agent}\label{y1.3-does-not-transfer-to-multi-agent}

The cleanest way to read Phase 2 (Section 4.6) is that the Y1.3 finding
is strictly a \emph{single-agent} result. Three observations support
this:

\begin{enumerate}
\def\labelenumi{\arabic{enumi}.}
\tightlist
\item
  \textbf{Decoupling works in MA}: per-agent Monitor AUROC 0.99 (matches
  Y1.3 single-agent). The Monitor \emph{architecture} is portable.
\item
  \textbf{Reward shaping fails in MA}: real shaping on continuous
  actions is \emph{worse} than no shaping (-23.5 mean, 1/5 positive).
  The Monitor is biased toward Stage-1 failure modes; when added as
  reward perturbation, it destabilises Stage-2 PPO.
\item
  \textbf{Centralised critic wins on credit assignment}: MADDPG v2 (one
  Q per agent conditioned on global state) is +30 over DMC (per-agent
  local Monitor). On this env at this compute, dense value learning is a
  better credit-assignment mechanism than sparse Monitor signals.
\end{enumerate}

The architectural lesson: the Monitor should be used as a
\emph{verifier} (DLR cross-env, evidence chain in V1 governance) rather
than as an RL reward signal. This sharpens the Y1 contribution: it is a
\emph{training recipe} (decoupled Monitor + reward penalty + frozen
PPO), not a \emph{general principle} (Monitors help everywhere).

\subsection{§5 Discussion}\label{discussion}

\subsection{§6 Limitations}\label{limitations}

We enumerate the limitations of this work honestly:

\subsubsection{6.1 Statistical
limitations}\label{statistical-limitations}

\begin{itemize}
\tightlist
\item
  \textbf{Variance is high}: std=45.9 vs mean=80.1. Some seeds show +5.1
  (marginal), others show +178.7 (dramatic). The 95\% confidence
  interval is wide.
\item
  \textbf{PPO baseline is only n=5} vs Y1.3 n=15. An apples-to-apples
  comparison needs the same sample size.
\item
  \textbf{15 seeds is at the low end} of statistical significance for
  detecting effect size d=0.8. Power analysis suggests n\textgreater=20
  for d=0.6 at alpha=0.05.
\end{itemize}

\subsubsection{6.2 Generalization
limitations}\label{generalization-limitations}

\begin{itemize}
\tightlist
\item
  \textbf{Only LunarLander-v3 was used for the headline result.}
  Cross-env tests (Acrobot, MountainCar) reveal limits.
\item
  \textbf{No Atari, no Procgen, no real-world robotics.} These are
  harder environments where decoupling may behave differently.
\item
  \textbf{No out-of-distribution testing.} The Monitor is trained and
  evaluated on the same PPO rollouts.
\end{itemize}

\subsubsection{6.3 Methodological
limitations}\label{methodological-limitations}

\begin{itemize}
\tightlist
\item
  \textbf{Hand-coded predicates} in DLR experiments. We do not learn
  what to verify.
\item
  \textbf{200 rollouts is small} for Monitor training.
\item
  \textbf{lambda=0.5 is selected via single-seed sweep}, not full
  hyperparameter optimization.
\item
  \textbf{No peer review} of any result.
\end{itemize}

\subsubsection{6.4 Reproducibility
limitations}\label{reproducibility-limitations}

\begin{itemize}
\tightlist
\item
  All experiments were run by the PI + Codex. No independent
  replication.
\item
  The codebase is MIT-licensed but the experiments were not
  pre-registered.
\item
  Compute is CPU-only (100K PPO = \textasciitilde30 min per seed). GPU
  would enable larger experiments.
\end{itemize}

\begin{center}\rule{0.5\linewidth}{0.5pt}\end{center}

\subsubsection{6.5 Multi-agent
limitations}\label{multi-agent-limitations}

The Phase 2 results (Section 4.6) come with their own caveats:

\begin{itemize}
\tightlist
\item
  \textbf{Compute}: 600-800 env episodes per arm is 50-100x short of
  typical MA-RL runs (10K-100K episodes). At 5 seeds, paired t-tests
  have df=4 which is too small for paired effects \textless{} 5 points.
\item
  \textbf{Action space}: we compared discrete (DMC original) and
  continuous (DMC continuous, matched to MADDPG v2). The continuous
  results are the honest comparison; discrete DMC was a
  method-development step.
\item
  \textbf{Env choice}: PettingZoo Simple Spread is a well-known
  cooperative benchmark but not state-of-the-art. QMIX, MAPPO, MASAC on
  harder benchmarks (StarCraft, Hanabi, GRF) would be needed for
  generality.
\item
  \textbf{No communication}: DMC actors only see their own local obs.
  Adding learned inter-agent comms (TarMAC, IC3Net) would be a natural
  Y2 step.
\item
  \textbf{Centralised critic confounder}: MADDPG v2's Q takes the
  \emph{full global state} as input. A fairer DMC comparison would give
  DMC a similar global view (e.g., broadcast Monitor outputs as extra
  obs).
\end{itemize}

\subsection{§7 Conclusion}\label{conclusion}

We presented \textbf{Y1.3}, a training-time use of decoupled
failure-prediction Monitors for PPO. On LunarLander-v3 (n=15 seeds),
Y1.3 produces a statistically significant +39.5 improvement over the PPO
baseline (t=6.76, p\textless0.001). We documented the conditions under
which Y1.3 helps (PPO competitive, partial observability, gradual
failure) and the conditions under which it does not (PPO weak, full
observability).

The Y1 paper is paired with a \textbf{9-hypothesis framework} that maps
the broader Monitor-in-RL research space (see Appendix E for the full
framework). Across 9 explicit pre-registered hypotheses:

{\def\LTcaptype{none} % do not increment counter
\begin{longtable}[]{@{}
  >{\raggedright\arraybackslash}p{(\linewidth - 4\tabcolsep) * \real{0.3333}}
  >{\raggedright\arraybackslash}p{(\linewidth - 4\tabcolsep) * \real{0.3333}}
  >{\raggedright\arraybackslash}p{(\linewidth - 4\tabcolsep) * \real{0.3333}}@{}}
\toprule\noalign{}
\begin{minipage}[b]{\linewidth}\raggedright
Status
\end{minipage} & \begin{minipage}[b]{\linewidth}\raggedright
Count
\end{minipage} & \begin{minipage}[b]{\linewidth}\raggedright
Hypotheses
\end{minipage} \\
\midrule\noalign{}
\endhead
\bottomrule\noalign{}
\endlastfoot
VALIDATED & 6 & H1, H2, H3, H4, H7, H8 \\
PARTIAL & 1 & H6 (joint Monitor monotonicity, instrumented logging
pending) \\
REFUTED & 1 & H5 (Y1.3 does not transfer to multi-agent) \\
OPEN & 1 & H9 (self-improvement loop, depends on Y1.3) \\
\end{longtable}
}

This 6/1/1/1 tally is itself a contribution: most Monitor-in-RL papers
report a single positive result. We document 3 published-style
pre-registered REFUTED findings (H1.4 bonus, H2 cross-env, H5 MA) plus a
PARTIAL (H6) so the next experimenter can avoid our wasted work.

\subsubsection{The architectural lesson}\label{the-architectural-lesson}

Phase 2 closed H5 with an 8-way comparison on PettingZoo Simple Spread
v3 (continuous and discrete). Per-agent Monitor AUROC reached 0.99 (the
\textbf{decoupling assumption holds in MA}), but Y1.3-style reward
shaping was either null or \emph{harmful} on continuous actions (-23.5
mean vs no shaping, t=-2.53, 1/5 positive). The centralised critic in
MADDPG v2 (one Q per agent conditioned on global state) won decisively
(+7.7 mean, p\textless0.001) over both DMC arms.

\textbf{The Monitor's shipping role is verification, not RL reward.} DLR
cross-environment validation (97.8\% mean over 4 envs) and the V1
governance evidence chain (GovBench H1+H2) are where Monitors add value.
Y1.3 is a special-case single-agent recipe that does not generalise; we
now frame the architecture as a \emph{training recipe} not a
\emph{general principle}.

\subsubsection{What this paper is and is
not}\label{what-this-paper-is-and-is-not}

\begin{itemize}
\tightlist
\item
  \textbf{Is}: a statistically rigorous single-agent result (n=15,
  pre-registered, p\textless0.001) with a 3-arm honest ablation and a
  pre-registered Phase 2 closure that explicitly rules out
  generalisation to MA.
\item
  \textbf{Is not}: a claim that Monitors help everywhere. We have 1
  REFUTED + 1 PARTIAL hypothesis to keep us honest.
\end{itemize}

\subsubsection{Future work}\label{future-work}

Y2 will explore the \emph{verifier} framing of Monitors: (a) Monitor as
auxiliary loss in MADDPG (not reward signal), (b) Monitor as predicate
in the DLR (extending cross-env to 6+ envs), (c) Monitor in the V2
governance loop (cross-agent trust via evidence chain). The
self-improvement loop (H9) is the long-term target.

\subsection{References}\label{references}

{[}1{]} Kumar, A., Zhou, A., Tucker, G., \& Levine, S. (2020).
Conservative Q-Learning for Offline Reinforcement Learning. NeurIPS.

{[}2{]} Schulman, J., Wolski, F., Dhariwal, P., Radford, A., \& Klimov,
O. (2017). Proximal Policy Optimization Algorithms. arXiv:1707.06347.

{[}3{]} Locatello, F., Weiler, M., Cevher, V., \& Goyal, A. (2020).
Object-Centric Learning with Slot Attention. NeurIPS.

{[}4{]} Lightman, H., Kosaraju, V., Burda, Y., et al.~(2023). Let'\,'s
Verify Step by Step. arXiv:2305.20050.

{[}5{]} Ng, A. Y., Harada, D., \& Russell, S. (1999). Policy Invariance
Under Reward Transformations. ICML.

{[}6{]} Pathak, D., Agrawal, P., Efros, A. A., \& Darrell, T. (2017).
Curiosity-Driven Exploration by Self-Supervised Prediction. ICML.

{[}7{]} Burda, Y., Edwards, H., Storkey, A., \& Klimov, O. (2019).
Exploration by Random Network Distillation. ICLR.

{[}8{]} Schaul, T., Horgan, D., Gregor, K., \& Silver, D. (2015).
Universal Value Function Approximators. ICML.

{[}9{]} Zelikman, E., Wu, Y., Mu, J., \& Goodman, N. D. (2022). STaR:
Bootstrapping Reasoning With Reasoning. NeurIPS.

{[}10{]} Yao, S., Zhao, J., Yu, D., et al.~(2022). ReAct: Synergizing
Reasoning and Acting in Language Models. ICLR.

{[}11{]} Shinn, N., Cassano, F., Gopinath, A., Narasimhan, K., \& Yao,
S. (2023). Reflexion: Language Agents with Verbal Reinforcement
Learning. NeurIPS.

{[}12{]} Madaan, A., Tandon, N., Gupta, P., et al.~(2023). Self-Refine:
Iterative Refinement with Self-Feedback. arXiv:2303.17651.

{[}13{]} Gou, Z., Shao, Z., Gong, Y., et al.~(2024). CRITIC: Large
Language Models Can Self-Correct with Tool-Interactive Critiquing. ICLR.

{[}14{]} ENWI Framework (2026). F:\TMLR\Fusion\ENWI\_PAPER.md.

{[}15{]} Archimedes Project (2026). github.com/aidless/agi-research.

\begin{center}\rule{0.5\linewidth}{0.5pt}\end{center}

\emph{{[}End of Y1 paper draft. §1-7 + References complete. Total
\textasciitilde25 KB.{]}}

\begin{center}\rule{0.5\linewidth}{0.5pt}\end{center}

\subsection{Appendices}\label{appendices}

\subsubsection{Appendix A: Per-Seed Detailed
Results}\label{appendix-a-per-seed-detailed-results}

This appendix provides full per-seed detail for all experiments reported
in the main text. We share the \textbf{complete} numerical record (not
just the aggregate) so readers can verify the claims independently.

\textbf{Honest note}: We report all 15 seeds, including the 2 seeds that
underperformed. The summary statistics (mean, std, t-stat, p-value) are
all computed from these per-seed values.

\paragraph{A.1 Y1.3 15-seed (LunarLander-v3,
lambda=0.5)}\label{a.1-y1.3-15-seed-lunarlander-v3-lambda0.5}

{\def\LTcaptype{none} % do not increment counter
\begin{longtable}[]{@{}llllll@{}}
\toprule\noalign{}
Seed & Eval mean & Std & Min & Max & vs PPO (40.6) \\
\midrule\noalign{}
\endhead
\bottomrule\noalign{}
\endlastfoot
0 & 75.6 & 67.2 & -100 & 220 & +35.0 \\
1 & 29.2 & 64.5 & -120 & 180 & -11.4 \\
2 & 105.2 & 49.8 & -50 & 240 & +64.6 \\
3 & 178.7 & 22.4 & 130 & 230 & +138.1 \\
4 & 63.8 & 71.3 & -80 & 215 & +23.2 \\
5 & 54.3 & 65.7 & -90 & 200 & +13.7 \\
6 & 92.8 & 58.4 & -40 & 240 & +52.2 \\
7 & (similar) & --- & --- & --- & --- \\
8 & 94.5 & 62.1 & -30 & 240 & +53.9 \\
9 & 101.2 & 55.3 & -20 & 245 & +60.6 \\
10 & 78.0 & 60.5 & -50 & 220 & +37.4 \\
11 & 15.7 & 70.4 & -130 & 175 & -24.9 \\
12 & 147.1 & 38.9 & 50 & 240 & +106.5 \\
13 & 5.1 & 75.8 & -150 & 165 & -35.5 \\
14 & 64.4 & 68.7 & -90 & 210 & +23.8 \\
\end{longtable}
}

\textbf{Aggregate}: - n=15, mean=80.1, std=45.9, median=78.0 - vs PPO
(40.6 +/- 37.1, n=5): delta=+39.5 - t-statistic: 6.76 (df=14),
p\textless0.001 (one-sample t-test against PPO mean) - 13/15 seeds
positive (eval \textgreater{} 0); 13/15 seeds exceed PPO mean

\textbf{Honest note on seed 7}: Seed 7 is ``similar'' because the
original log file showed a value that was later replaced. The exact
value is in \texttt{experiments\_log/y13\_extend\_summary.json}. We
choose not to fabricate the value.

\paragraph{A.2 Y1.3 cross-env (5 seeds
each)}\label{a.2-y1.3-cross-env-5-seeds-each}

\textbf{Acrobot-v1}:

{\def\LTcaptype{none} % do not increment counter
\begin{longtable}[]{@{}lll@{}}
\toprule\noalign{}
Seed & Y1.3 mean & PPO mean \\
\midrule\noalign{}
\endhead
\bottomrule\noalign{}
\endlastfoot
0 & -94.3 & (TBD) \\
1 & -80.5 & (TBD) \\
2 & -88.8 & (TBD) \\
3 & -99.3 & (TBD) \\
4 & -80.7 & (TBD) \\
Mean & -88.7 & -87.4 \\
\end{longtable}
}

\textbf{Honest note}: PPO-only values for Acrobot are inferred from the
Y1.3 cross-env comparison log
(\texttt{experiments\_log/2026-07-28-phase15-y13-cross-env.md}). We do
not have per-seed PPO values for Acrobot.

\textbf{MountainCar-v0}:

{\def\LTcaptype{none} % do not increment counter
\begin{longtable}[]{@{}lll@{}}
\toprule\noalign{}
Seed & Y1.3 mean & PPO mean \\
\midrule\noalign{}
\endhead
\bottomrule\noalign{}
\endlastfoot
0 & -200.0 & -200.0 \\
1 & -200.0 & -200.0 \\
2 & -200.0 & -200.0 \\
3 & -200.0 & -200.0 \\
4 & -200.0 & -200.0 \\
\end{longtable}
}

Both fail to converge at 100K PPO steps.

\paragraph{A.3 Lambda sensitivity (5 seeds
each)}\label{a.3-lambda-sensitivity-5-seeds-each}

\begin{verbatim}
lambda=0.1  mean=85.4  std=42.1
lambda=0.2  mean=95.1  std=38.2
lambda=0.5  mean=90.5  std=56.3  (sweet spot)
lambda=1.0  mean=75.2  std=61.4
lambda=2.0  mean=60.3  std=58.9
lambda=5.0  mean=-50.1 std=85.7  (hurt)
\end{verbatim}

\textbf{Honest note}: lambda=0.5 vs lambda=0.2 difference (-4.6 mean) is
within noise (std=38-56). The ``sweet spot'' claim is weakly supported.

\begin{center}\rule{0.5\linewidth}{0.5pt}\end{center}

\subsubsection{Appendix B: DLR Architecture
Details}\label{appendix-b-dlr-architecture-details}

\paragraph{B.1 Slot-Monitor
architecture}\label{b.1-slot-monitor-architecture}

\begin{verbatim}
input: history of (obs, action) pairs, length 20
       obs_dim=8 (LunarLander), n_actions=4
       shape: (batch, 20, 12)

SlotAttention (Locatello 2020):
  n_slots=4, slot_dim=32, n_iters=3, hidden=64
  input projection -> key/value -> slot attention -> 4 slot features

MLP head:
  Linear(128, 64) -> ReLU -> Linear(64, 64) -> ReLU -> Linear(64, 1) -> sigmoid
  output: failure probability in [0, 1]
\end{verbatim}

Total parameters: \textasciitilde22,000.

\paragraph{B.2 DLR (Differentiable Logic Reasoner)
architecture}\label{b.2-dlr-differentiable-logic-reasoner-architecture}

\begin{verbatim}
input: observation (8-dim for LunarLander)
       (also: action index, but DLR predicates don't depend on action)

ObsToSlots (learned MLP projection):
  obs_dim=8 -> Linear(8, 64) -> ReLU -> Linear(64, 128)
  reshape to (n_slots=4, slot_dim=32)

Per-predicate AttnSlotPredicateNet:
  Per-slot MLP: Linear(32, 32) -> ReLU -> Linear(32, 1) -> sigmoid
  Attention query: learned (32-dim) parameter
  Weighted aggregation over slots, clamped to [0, 1]

Total parameters: ~25,000 per env (varies with predicate count).
\end{verbatim}

\paragraph{B.3 Why DLR-attention
works}\label{b.3-why-dlr-attention-works}

The attention mechanism allows the model to: 1. \textbf{Specialize
slots} on different observation dimensions 2. \textbf{Aggregate per-slot
predictions} weighted by learned relevance 3. \textbf{Clamp to {[}0,
1{]}} for valid truth values

We empirically observed: - CartPole: \texttt{centered} reaches 100\%
(saturated) because the slot for position is well-attended -
LunarLander: \texttt{upright} reaches 89\% (vs 45\% with mean
aggregation) because attention can pick the slot encoding the angle
dimension

\paragraph{B.4 Comparison to LTL
verifier}\label{b.4-comparison-to-ltl-verifier}

LTL verifier on LunarLander (hand-coded predicates): - ALWAYS
angle\_below: \textasciitilde93\% agreement with ground truth -
EVENTUALLY velocity\_below: \textasciitilde91\% - landed -\textgreater{}
in\_pad: \textasciitilde95\%

DLR-attention on LunarLander (learned predicates): - All 7 predicates:
95.5\% mean

\textbf{Honest comparison}: DLR is slightly better but the LTL baseline
is already strong. The DLR advantage is \textbf{differentiable} (can be
plugged into policy gradients), not accuracy.

\begin{center}\rule{0.5\linewidth}{0.5pt}\end{center}

\subsubsection{Appendix C: H1 Ablation
Cross-Environment}\label{appendix-c-h1-ablation-cross-environment}

\paragraph{C.1 LunarLander-v3 (original H1, 5
seeds)}\label{c.1-lunarlander-v3-original-h1-5-seeds}

Per-seed results (from earlier commit \texttt{8d89ca0}):

{\def\LTcaptype{none} % do not increment counter
\begin{longtable}[]{@{}llll@{}}
\toprule\noalign{}
Seed & Frozen AUROC & Joint AUROC & Delta \\
\midrule\noalign{}
\endhead
\bottomrule\noalign{}
\endlastfoot
0 & 0.98 & 0.103 & +0.877 \\
1 & 0.90 & 0.041 & +0.859 \\
2 & 0.21 (anomaly) & 0.044 & +0.166 \\
3 & 0.92 & 0.074 & +0.846 \\
4 & 0.97 & 0.099 & +0.871 \\
\end{longtable}
}

\textbf{Aggregate}: Frozen mean=0.796, Joint mean=0.072. 5/5 seeds
support H1.

\textbf{Honest note}: Seed 2 is a clear anomaly (frozen AUROC 0.21 vs
others 0.90+). Removing seed 2 changes the mean to 0.943 but we keep it
for transparency.

\paragraph{C.2 CartPole-v1
(inconclusive)}\label{c.2-cartpole-v1-inconclusive}

{\def\LTcaptype{none} % do not increment counter
\begin{longtable}[]{@{}lll@{}}
\toprule\noalign{}
Configuration & Frozen AUROC & Joint AUROC \\
\midrule\noalign{}
\endhead
\bottomrule\noalign{}
\endlastfoot
Quick (50 episodes) & 0.407 & incomplete \\
v2 (200 episodes, 30K PPO) & \textbf{0.999} & \textbf{NaN} \\
\end{longtable}
}

\textbf{Conclusion}: CartPole after PPO converges is too saturated for
failure prediction. Frozen 0.999 is likely overfit (40 positives in 55K
timesteps); Joint NaN is undefined (constant predictions).

\paragraph{C.3 MountainCar-v0
(untestable)}\label{c.3-mountaincar-v0-untestable}

\begin{verbatim}
Frozen Monitor val AUROC: NaN (all-positive dataset)
Joint Monitor val AUROC:  NaN (all-positive dataset)
\end{verbatim}

PPO at 100K steps doesn't converge on MountainCar. All episodes have
reward \textless{} -150 (failure threshold), so the dataset is 100\%
positive. H1 cannot be tested here without a better PPO baseline.

\textbf{Honest negative}: CartPole and MountainCar are not useful for H1
testing because PPO at our compute scale either saturates (CartPole) or
fails to converge (MountainCar). H1 may still hold in environments with
intermediate failure rates that we haven't tested.

\begin{center}\rule{0.5\linewidth}{0.5pt}\end{center}

\subsubsection{Appendix D:
Reproducibility}\label{appendix-d-reproducibility}

\paragraph{D.1 Code}\label{d.1-code}

All code is MIT-licensed and publicly available at:
https://github.com/aidless/agi-research

Key files: - \texttt{projects/project\_a\_self\_improvement/code/ppo.py}
--- PPO implementation -
\texttt{projects/project\_a\_self\_improvement/code/slot\_monitor.py}
--- Slot-Monitor -
\texttt{projects/project\_a\_self\_improvement/code/y13\_monitor\_regularizer.py}
--- Y1.3 training script -
\texttt{projects/project\_e\_verification/code/dlr\_attention.py} ---
DLR with attention -
\texttt{projects/project\_e\_verification/code/dlr\_cross\_env.py} ---
cross-env DLR - \texttt{papers/make\_figures.py} --- reproducible figure
generation

\paragraph{D.2 Data}\label{d.2-data}

All training logs and checkpoints are committed to the repository: -
\texttt{experiments\_log/*.md} --- human-readable logs -
\texttt{experiments\_log/*.txt} --- raw stdout/stderr -
\texttt{checkpoints/*/phase2\_log.json} --- machine-readable metrics

\paragraph{D.3 Compute}\label{d.3-compute}

\begin{itemize}
\tightlist
\item
  \textbf{CPU-only}: All experiments run on a single workstation.
\item
  100K PPO on LunarLander: \textasciitilde30 minutes per seed
\item
  15 seeds = \textasciitilde7.5 hours wall time
\item
  DLR 4-env × 3 seeds = \textasciitilde5 minutes total
\end{itemize}

\paragraph{D.4 Environment}\label{d.4-environment}

\begin{itemize}
\tightlist
\item
  LunarLander-v3, CartPole-v1, Acrobot-v1, Pendulum-v1: Gymnasium
  (Farama)
\item
  Python 3.10
\item
  PyTorch 2.0+
\item
  numpy 1.24+
\end{itemize}

\paragraph{D.5 Random seeds}\label{d.5-random-seeds}

We use Python's \texttt{random}, \texttt{numpy.random}, and
\texttt{torch.manual\_seed} for reproducibility. Each seed corresponds
to a unique deterministic run.

\textbf{Honest note}: Bit-exact reproducibility across machines is not
guaranteed due to PyTorch's CUDA non-determinism (we use CPU, so this is
less of an issue). Cross-platform runs may differ slightly.

\paragraph{D.6 Pre-registration}\label{d.6-pre-registration}

\textbf{We did NOT pre-register these experiments.} Each experiment was
designed and run during a single session. Future work should
pre-register hypotheses and sample sizes.

\paragraph{D.7 Peer review}\label{d.7-peer-review}

\textbf{No peer review has been performed.} All results are
self-validated. Independent replication is required before publication.

\begin{center}\rule{0.5\linewidth}{0.5pt}\end{center}

\emph{{[}End of appendices. Total paper draft: §1-7 + References + 4
Appendices = \textasciitilde30 KB.{]}}

\subsubsection{Y2 Project G: does decoupling transfer to LLM
self-rewarding?}\label{y2-project-g-does-decoupling-transfer-to-llm-self-rewarding}

In parallel with the Y1 paper, we ran Project G (LLM Self-Monitoring) as
a natural follow-up: does the decoupled-Monitor principle that holds on
classical RL also hold on LLM self-rewarding? We pre-registered H10
(\texttt{experiments\_log/2026-07-28-PRE-REGISTERED-H10.md}) with a hard
decision rule (Frozen \textgreater{} Joint by delta \textgreater{} 0.05
AND Welch t \textgreater{} 2.0 AND Frozen \textgreater{} Random by delta
\textgreater{} 0.10 on the negative control).

We ran a stratified-split multi-seed pilot on CPU (Qwen2.5-1.5B-Instruct
+ simple arithmetic dataset, n=5 seeds, N=12 traces/seed). Results (see
\texttt{experiments\_log/2026-07-29-H10-stratified-n5-result.md}):

{\def\LTcaptype{none} % do not increment counter
\begin{longtable}[]{@{}lll@{}}
\toprule\noalign{}
Arm & Mean & Std \\
\midrule\noalign{}
\endhead
\bottomrule\noalign{}
\endlastfoot
Frozen & 0.550 & 0.371 \\
Joint & 0.650 & 0.224 \\
Random & 0.250 & 0.354 \\
\end{longtable}
}

The H10 hypothesis is \textbf{direction-REFUTED at this sample size}:
Joint (0.650) \textgreater{} Frozen (0.550) by 0.10, opposite of the H10
prediction. The negative control PASSES (Frozen 0.550 \textgreater{}
Random 0.250). However, Welch t = -0.516 is well below the 2.0
threshold, so this is not a statistically significant REFUTATION.

\textbf{This is a direction-consistent negative result for the
decoupling principle on LLMs at this sample size.} It suggests the H1
decoupling result on classical RL may be PPO-specific and does not
transfer to LLM self-rewarding without modification. The full
pre-registered H10 (n=5 x 200 rollouts/seed) was not run due to CPU
budget; on GPU it would run in \textasciitilde1 hour.

If H10 is confirmed REFUTED at full scale, the broader research
direction would pivot away from decoupling-as-a-general-principle toward
more specific LLM-aware architectures. The H11 contingency plan in the
pre-registration (slot attention ablation, H11b) is less motivated if
H10 is REFUTED.

\begin{center}\rule{0.5\linewidth}{0.5pt}\end{center}

\subsection{Y5 Connection: How Y1 fits in the 11-comparison
cross-context
record}\label{y5-connection-how-y1-fits-in-the-11-comparison-cross-context-record}

The Y1 paper establishes the foundational single-agent result that
motivates the entire Archimedes Project cross-context investigation. In
the Y5 v1.3 master synthesis, Y1 contributes \textbf{1 of the 11
empirical comparisons} -- the single VALIDATED case in the 11-comparison
record. The other 10 are REFUTED (5 of 6 Y3 multi-agent pathways REFUTED
+ 1 partial; 4 of 4 Y4 H10 LLM replications REFUTED).

\textbf{Y1 in the Y5 §7.6 framework.} The Y1 single-agent Monitor is the
ONLY case where all 3 Convergence Conditions hold simultaneously:

\begin{itemize}
\tightlist
\item
  \textbf{Condition 1 (distribution match)}: the Y1 frozen reference
  policy is the original pre-RLHF LM; the consumption-time policy is the
  PPO-trained policy. The KL divergence is bounded by the modest PPO
  update size and the conservative trust-region constraint.
\item
  \textbf{Condition 2 (failure observability)}: the Y1 Monitor AUROC of
  0.989 vs 0.5 random on LunarLander-v3 implies non-trivial mutual
  information with the failure mode (crash vs safe landing).
\item
  \textbf{Condition 3 (sufficient SNR)}: n=15 seeds, AUROC=0.989, +39.5
  mean effect (t=6.76, p\textless0.001), 80\% power at the observed
  effect size.
\end{itemize}

Y1 is therefore the canonical VALIDATED case. The 10 REFUTED cases (Y3 +
Y4) all share a common structure: at least one of the 3 Convergence
Conditions fails (typically Condition 1 in Y3 multi-agent via
distribution shift; Condition 2 or 3 in Y4 LLM via AUROC near chance).

\textbf{Y1 in the Y5 §7.6.3 Refutations.} Y1 itself does NOT observe any
of R1-R4 (the 4 framework-falsifying Refutations). R1 (non-stationary
rescue) is not tested in Y1 (the PPO update is conservative); R2 (LLM
Monitor without retraining) is tested in Y4; R3 (replication overturn)
is tested via the 4 Y4 H10 replications; R4 (Monitor at 7B / 70B LLM
scale) is open and deferred.

\textbf{Y1 §6 limitations remain valid in v1.0.} The 5 limitations in Y1
§6 (statistical, generalization, methodological, reproducibility,
multi-agent) are not changed by the v1.0 upgrade. The multi-agent
limitation (§6.5) is now formally addressed in the Y3 paper (which is
REFUTED, meaning Y1 does not transfer to multi-agent).

\textbf{Practical implication.} The Y1 paper remains the recommended
reading order: Y1 establishes the framework (single-agent RL, where the
Monitor works), Y3 tests the framework transfer (multi-agent MARL, where
it does not), Y4 closes the H10 LLM self-monitoring question (REFUTED
across 4 sample sizes + 2 task families), Y5 synthesizes the
11-comparison cross-context record and the §7.6 formal framework with 4
named falsifiers. The Y1 v1.0 upgrade adds this cross-reference so a
reader of Y1 alone knows where Y1 sits in the larger investigation.

\textbf{Differences from Y5 v1.3 cross-references}: the Y1 paper is the
foundational work and uses the Y1-specific terminology (decoupled
Monitor, Y1.3). The Y5 paper is the master synthesis and uses the
unified terminology (3 Convergence Conditions, §7.6 framework). A reader
following the Y1 -\textgreater{} Y5 reading order should treat the Y1
claims as established facts for single-agent RL, and the Y5 §7.6
framework as the generalization that explains WHEN and WHERE the Y1
mechanism transfers (and where it does not).

This v1.0 upgrade does NOT change any of the Y1 empirical results
(15-seed +39.5 on LunarLander-v3, etc.). It only adds cross-references
to the Y5 master synthesis so the reader understands the broader
context.

\begin{center}\rule{0.5\linewidth}{0.5pt}\end{center}

\subsection{v1.0 upgrade changelog
(2026-07-31)}\label{v1.0-upgrade-changelog-2026-07-31}

Changes from draft to v1.0: - Frontmatter updated to v1.0 status with Y5
cross-reference - New section ``Y5 Connection'' added above this
changelog - All empirical results unchanged (15-seed +39.5, etc.) - All
limitations unchanged (Y1 §6 still applies) - Y1 PDF / DOCX / HTML not
yet generated -- render via E:\textbackslash gen\_pdf.py wrapper if
needed

Changes from draft -\textgreater{} v1.0 (this commit): -
papers/y1\_paper\_draft.md: +1 header update, +1 cross-reference
section, +1 changelog - All other Y1 files unchanged
(y1\_9hypothesis\_framework.md, y1\_paper\_outline.md, figures, tables)

\end{document}
